\subsection{Approximate Factorization Method}
\label{subsec:apf}

The approximate factorization (APF) method advances the
semi-discrete energy equation in pseudo-time by splitting the
two-dimensional update into two successive one-dimensional
linear systems. This approach avoids solving a single
multidimensional system at every iteration. Instead, the
resulting systems are tridiagonal and can be solved efficiently
using the Thomas algorithm.

\subsubsection{Approximate Factorization Formulation}

Let the temperature correction during one pseudo-time step be
defined as
\begin{equation}
\delta T_{ij}
=
T_{ij}^{n+1}-T_{ij}^{n},
\label{eq:apf_temperature_correction}
\end{equation}
where $n$ and $n+1$ denote the current and next pseudo-time
levels, respectively.

The semi-discrete energy equation is written as
\begin{equation}
\frac{\mathrm{d}T_{ij}}{\mathrm{d}t}
=
\mathcal{L}_x(T)_{ij}
+
\mathcal{L}_y(T)_{ij}
+
S_{ij}
=
\mathcal{R}_{ij}(T),
\label{eq:apf_residual}
\end{equation}
where $\mathcal{L}_x$ and $\mathcal{L}_y$ represent the
streamwise and wall-normal transport operators, respectively,
and $S_{ij}$ is the viscous-dissipation source term. The
spatial derivatives are approximated using the central
finite-difference expressions introduced in the previous section.

For compactness, the two spatial operators are expressed as
\begin{align}
\mathcal{L}_x(T)_{ij}
={}&
-\frac{
\alpha_c
\left(
u_{i+1,j}T_{i+1,j}
-u_{i-1,j}T_{i-1,j}
\right)
}{2\Delta x}
\nonumber\\
&+
\frac{1}{\mathrm{Pe}}
\frac{
T_{i+1,j}-2T_{ij}+T_{i-1,j}
}{\Delta x^2},
\label{eq:apf_x_operator}
\\[0.25cm]
\mathcal{L}_y(T)_{ij}
={}&
-\frac{
\alpha_c
\left(
v_{i,j+1}T_{i,j+1}
-v_{i,j-1}T_{i,j-1}
\right)
}{2\Delta y}
\nonumber\\
&+
\frac{1}{\mathrm{Pe}}
\frac{
T_{i,j+1}-2T_{ij}+T_{i,j-1}
}{\Delta y^2},
\label{eq:apf_y_operator}
\end{align}
where
\begin{equation}
\mathrm{Pe}=\mathrm{Re}\,\mathrm{Pr},
\qquad
\mathrm{Br}=\mathrm{Ec}\,\mathrm{Pr},
\end{equation}
and $\alpha_c$ is the convection coefficient used in the
implementation, with default value $\alpha_c=1$.

The viscous-dissipation source term is evaluated from the
prescribed velocity field:
\begin{equation}
\boxed{
S_{ij}
=
\frac{\mathrm{Br}}{\mathrm{Pe}}
\left[
2\left(\frac{\partial u}{\partial x}\right)^2
+
2\left(\frac{\partial v}{\partial y}\right)^2
+
\left(
\frac{\partial v}{\partial x}
+
\frac{\partial u}{\partial y}
\right)^2
\right]_{ij}.
}
\label{eq:apf_viscous_source}
\end{equation}
The velocity gradients are approximated using central differences.
Since $\mathrm{Br}/\mathrm{Pe}=\mathrm{Ec}/\mathrm{Re}$,
Equation~\eqref{eq:apf_viscous_source} is consistent with the
source term in the governing energy equation.

The residual at the current pseudo-time level is therefore
\begin{equation}
\boxed{
\mathcal{R}_{ij}^{n}
=
\mathcal{L}_x(T^n)_{ij}
+
\mathcal{L}_y(T^n)_{ij}
+
S_{ij}.
}
\label{eq:apf_current_residual}
\end{equation}

The APF formulation approximates the multidimensional
correction equation through the factorized system
\begin{equation}
\boxed{
\left(I-\Delta t\,\mathcal{L}_x\right)
\left(I-\Delta t\,\mathcal{L}_y\right)
\delta T
=
\Delta t\,\mathcal{R}^{n},
}
\label{eq:apf_factorized_equation}
\end{equation}
where $I$ denotes the identity operator and $\Delta t$ is
the pseudo-time step. The product of the two operators
introduces a cross term proportional to $\Delta t^2$.
Consequently, the factorized equation approximates the
corresponding unfactorized multidimensional correction equation.

Introducing an intermediate correction
$\delta\widetilde{T}$, the factorized equation is solved in
two stages:
\begin{align}
\left(I-\Delta t\,\mathcal{L}_x\right)
\delta\widetilde{T}
&=
\Delta t\,\mathcal{R}^{n},
\label{eq:apf_x_stage}
\\[0.2cm]
\left(I-\Delta t\,\mathcal{L}_y\right)
\delta T
&=
\delta\widetilde{T}.
\label{eq:apf_y_stage}
\end{align}

The first stage solves for the intermediate correction along
the streamwise direction. The second stage uses this
intermediate field to obtain the final correction along the
wall-normal direction. The temperature is then updated as
\begin{equation}
\boxed{
T_{ij}^{n+1}
=
T_{ij}^{n}+\delta T_{ij}.
}
\label{eq:apf_temperature_update}
\end{equation}

\subsubsection{Construction of the Spatial Residual}

The right-hand side of the first factorized system is
constructed using the function \texttt{build\_rhs}. For each
physical wall-normal index $j$, this function evaluates the
streamwise and wall-normal convective and diffusive
contributions, calculates the viscous-dissipation source,
and assembles the residual vector.

Defining the transport contribution as
\begin{equation}
F_{ij}
=
\mathcal{L}_x(T)_{ij}
+
\mathcal{L}_y(T)_{ij},
\label{eq:apf_transport_contribution}
\end{equation}
the residual is
\begin{equation}
\mathcal{R}_{ij}(T)=F_{ij}+S_{ij}.
\label{eq:apf_residual_assembly}
\end{equation}

The source term is calculated using Equation~\eqref{eq:apf_viscous_source}. The resulting algorithm corresponds to the implementation
of \texttt{build\_rhs}.

\begin{algorithm}[H]
\caption{Construction of the Spatial Residual Vector}
\label{alg:apf_build_rhs}
\begin{algorithmic}[1]
\Require Mesh dimensions $n_x,n_y$, wall-normal index $j$,
temperature field $T$, velocity fields $u,v$, parameters
$\mathrm{Pe}$ and $\mathrm{Br}$, domain dimensions $L_x,L_y$,
and convection coefficient $\alpha_c$
\Ensure Residual vector $\mathbf{r}^{(j)}$
\State $\Delta x \gets L_x/n_x$
\State $\Delta y \gets L_y/n_y$
\State Initialize $\mathbf{r}^{(j)}$ as a zero vector of
length $n_x+2$
\For{$i=1,\ldots,n_x$}
    \State Evaluate the transport contribution
    $F_{ij}$ using Equation~\eqref{eq:apf_transport_contribution}
    \State Calculate the four velocity gradients using
    central differences
    \State Evaluate the source term $S_{ij}$ using
    Equation~\eqref{eq:apf_viscous_source}
    \State $\mathbf{r}^{(j)}_i \gets F_{ij}+S_{ij}$
\EndFor
\State \Return $\mathbf{r}^{(j)}$
\end{algorithmic}
\end{algorithm}

The returned vector contains residual values at physical
streamwise indices; its ghost-cell entries remain zero.
Accordingly, the right-hand side of the first factorized
system is $\Delta t\,\mathbf{r}^{(j)}$.

\subsubsection{Boundary Conditions in the Factorized Systems}

The boundary conditions are imposed through the boundary rows
of the factorized systems. In terms of temperature corrections,
the four boundary conditions are

\begin{align}
\text{Inlet:}\qquad
&\delta T_{0,j}+\delta T_{1,j}=0,
\\
\text{Outlet:}\qquad
&\delta T_{n_x,j}-\delta T_{n_x+1,j}=0,
\\
\text{Bottom wall:}\qquad
&\delta T_{i,0}+\delta T_{i,1}=0,
\\
\text{Top wall:}\qquad
&\delta T_{i,n_y}+\delta T_{i,n_y+1}=0.
\end{align}

The inlet and outlet conditions are incorporated into the
streamwise system, while the wall conditions are imposed in the
wall-normal system. The corresponding matrix rows contain two
nonzero coefficients: $(1,1)$ for prescribed-temperature
boundaries and $(1,-1)$ for the outlet zero-gradient condition,
with the latter ordered as in the matrix implementation.

\subsubsection{Streamwise System}

For each fixed wall-normal index $j$, the first APF stage
requires the solution of
Equation~\eqref{eq:apf_x_stage} along the streamwise direction.
For a physical cell $(i,j)$, the resulting equation is
\begin{equation}
\boxed{
a^x_{W,i}\,
\delta\widetilde{T}_{i-1,j}
+
a^x_P\,
\delta\widetilde{T}_{ij}
+
a^x_{E,i}\,
\delta\widetilde{T}_{i+1,j}
=
\Delta t\,\mathcal{R}_{ij}^{n},
}
\label{eq:apf_x_tridiagonal}
\end{equation}
where the coefficients are
\begin{align}
a^x_{W,i}
&=
-\frac{\alpha_c u_{i-1,j}\Delta t}{2\Delta x}
-\frac{\Delta t}{\mathrm{Pe}\,\Delta x^2},
\label{eq:apf_x_west}
\\[0.15cm]
a^x_P
&=
1+\frac{2\Delta t}{\mathrm{Pe}\,\Delta x^2},
\label{eq:apf_x_center}
\\[0.15cm]
a^x_{E,i}
&=
\frac{\alpha_c u_{i+1,j}\Delta t}{2\Delta x}
-\frac{\Delta t}{\mathrm{Pe}\,\Delta x^2}.
\label{eq:apf_x_east}
\end{align}

For each $j$, the unknown corrections are arranged in a
vector containing the physical cells and the two streamwise
ghost cells. The coefficient matrix is tridiagonal. The
right-hand side contains the residual values at the physical
cells, scaled by $\Delta t$, with zero entries at the ghost
indices.

The corresponding matrix is assembled by
\texttt{build\_x\_matrix}, as summarized below.

\begin{algorithm}[H]
\caption{Construction of the Streamwise Coefficient Matrix}
\label{alg:apf_build_x_matrix}
\begin{algorithmic}[1]
\Require Mesh dimensions $n_x,n_y$, wall-normal index $j$,
time step $\Delta t$, Peclet number $\mathrm{Pe}$, velocity
field $u$, domain length $L_x$, and convection coefficient
$\alpha_c$
\Ensure Streamwise matrix $\mathbf{A}_x^{(j)}$
\State $\Delta x \gets L_x/n_x$
\State Initialize $\mathbf{A}_x^{(j)}$ as a zero matrix of
size $(n_x+2)\times(n_x+2)$
\State Set the inlet row:
$A^x_{0,0}\gets1$, $A^x_{0,1}\gets1$
\State Set the outlet row:
$A^x_{n_x+1,n_x}\gets1$,
$A^x_{n_x+1,n_x+1}\gets-1$
\For{$i=1,\ldots,n_x$}
    \State $A^x_{i,i}
    \gets 1+\dfrac{2\Delta t}{\mathrm{Pe}\,\Delta x^2}$
    \State $A^x_{i,i+1}
    \gets \dfrac{\alpha_c\Delta t\,u_{i+1,j}}{2\Delta x}
    -\dfrac{\Delta t}{\mathrm{Pe}\,\Delta x^2}$
    \State $A^x_{i,i-1}
    \gets -\dfrac{\alpha_c\Delta t\,u_{i-1,j}}{2\Delta x}
    -\dfrac{\Delta t}{\mathrm{Pe}\,\Delta x^2}$
\EndFor
\State \Return $\mathbf{A}_x^{(j)}$
\end{algorithmic}
\end{algorithm}

For each physical wall-normal index, the resulting linear
system is
\begin{equation}
\boxed{
\mathbf{A}_x^{(j)}
\delta\widetilde{\mathbf{T}}^{(j)}
=
\Delta t\,\mathbf{r}^{(j)}.
}
\label{eq:apf_x_matrix_system}
\end{equation}
It is solved using the Thomas algorithm, producing one row
of the intermediate correction field.

\subsubsection{Wall-Normal System}

After the intermediate correction has been obtained, the
second APF stage solves
Equation~\eqref{eq:apf_y_stage} along the wall-normal
direction. For a physical cell $(i,j)$, the corresponding
equation is
\begin{equation}
\boxed{
a^y_{S,j}\,\delta T_{i,j-1}
+
a^y_P\,\delta T_{ij}
+
a^y_{N,j}\,\delta T_{i,j+1}
=
\delta\widetilde{T}_{ij}.
}
\label{eq:apf_y_tridiagonal}
\end{equation}
The coefficients are
\begin{align}
a^y_{S,j}
&=
-\frac{\alpha_c v_{i,j-1}\Delta t}{2\Delta y}
-\frac{\Delta t}{\mathrm{Pe}\,\Delta y^2},
\label{eq:apf_y_south}
\\[0.15cm]
a^y_P
&=
1+\frac{2\Delta t}{\mathrm{Pe}\,\Delta y^2},
\label{eq:apf_y_center}
\\[0.15cm]
a^y_{N,j}
&=
\frac{\alpha_c v_{i,j+1}\Delta t}{2\Delta y}
-\frac{\Delta t}{\mathrm{Pe}\,\Delta y^2}.
\label{eq:apf_y_north}
\end{align}

For each fixed streamwise index $i$, the unknown corrections
are arranged in a vector containing the physical cells and
the two wall-normal ghost cells. The right-hand side is
obtained from the corresponding column of the intermediate
correction field, with zero right-hand-side entries for the
boundary-condition rows.

The wall-normal coefficient matrix is assembled using
\texttt{build\_y\_matrix}.

\begin{algorithm}[H]
\caption{Construction of the Wall-Normal Coefficient Matrix}
\label{alg:apf_build_y_matrix}
\begin{algorithmic}[1]
\Require Mesh dimensions $n_x,n_y$, streamwise index $i$,
time step $\Delta t$, Peclet number $\mathrm{Pe}$, velocity
field $v$, domain length $L_y$, and convection coefficient
$\alpha_c$
\Ensure Wall-normal matrix $\mathbf{A}_y^{(i)}$
\State $\Delta y \gets L_y/n_y$
\State Initialize $\mathbf{A}_y^{(i)}$ as a zero matrix of
size $(n_y+2)\times(n_y+2)$
\State Set the bottom-wall row:
$A^y_{0,0}\gets1$, $A^y_{0,1}\gets1$
\State Set the top-wall row:
$A^y_{n_y+1,n_y}\gets1$,
$A^y_{n_y+1,n_y+1}\gets1$
\For{$j=1,\ldots,n_y$}
    \State $A^y_{j,j}
    \gets 1+\dfrac{2\Delta t}{\mathrm{Pe}\,\Delta y^2}$
    \State $A^y_{j,j+1}
    \gets \dfrac{\alpha_c\Delta t\,v_{i,j+1}}{2\Delta y}
    -\dfrac{\Delta t}{\mathrm{Pe}\,\Delta y^2}$
    \State $A^y_{j,j-1}
    \gets -\dfrac{\alpha_c\Delta t\,v_{i,j-1}}{2\Delta y}
    -\dfrac{\Delta t}{\mathrm{Pe}\,\Delta y^2}$
\EndFor
\State \Return $\mathbf{A}_y^{(i)}$
\end{algorithmic}
\end{algorithm}

For each physical streamwise index, the second stage is
written in matrix form as
\begin{equation}
\boxed{
\mathbf{A}_y^{(i)}
\delta\mathbf{T}^{(i)}
=
\delta\widetilde{\mathbf{T}}^{(i)}.
}
\label{eq:apf_y_matrix_system}
\end{equation}
This system is also tridiagonal and is solved using the
Thomas algorithm. Repeating the procedure over all physical
streamwise indices yields the final temperature correction
field.

In the present problem, the prescribed velocity field has
$v=0$. Consequently, the wall-normal convective terms in
the residual and wall-normal matrix coefficients vanish,
although the wall-normal diffusion terms remain.

\subsubsection{APF Iteration and Steady-State Convergence}

The complete APF procedure consists of repeatedly evaluating
the spatial residual, solving the streamwise and wall-normal
systems, updating the temperature field, and checking the
steady-state convergence criterion.

The RMS residual is defined over physical cells as
\begin{equation}
\boxed{
R_{\mathrm{RMS}}
=
\sqrt{
\frac{1}{n_xn_y}
\sum_{i=1}^{n_x}
\sum_{j=1}^{n_y}
\left[\mathcal{R}_{ij}(T)\right]^2
}.
}
\label{eq:apf_convergence_residual}
\end{equation}
It is evaluated using the same spatial residual as
Equation~\eqref{eq:apf_current_residual}. The iteration is
terminated when
\begin{equation}
\boxed{
R_{\mathrm{RMS}}<10^{-12}.
}
\label{eq:apf_convergence_criterion}
\end{equation}

\begin{algorithm}[H]
\caption{APF Pseudo-Time Marching to Steady State}
\label{alg:apf_complete_solver}
\begin{algorithmic}[1]
\Require Initial temperature field $T^0$, velocity fields
$u,v$, mesh dimensions $n_x,n_y$, time step $\Delta t$,
parameters $\mathrm{Pe}$ and $\mathrm{Br}$, domain dimensions
$L_x,L_y$, convection coefficient $\alpha_c$, tolerance
$\varepsilon=10^{-12}$, maximum iterations $N_{\max}$,
and snapshot interval $N_s$
\Ensure Final temperature field $T$, iteration and residual
histories, computation time, and temperature snapshots

\State $T\gets T^0$
\State Initialize iteration and residual histories as empty
\State Store $T^0$ as the initial snapshot at iteration $0$
\State Start the computation timer
\State $\mathrm{converged}\gets\mathrm{false}$

\For{$n=1,\ldots,N_{\max}$}
    \State Apply the temperature boundary conditions
    \State Initialize the intermediate correction field
    $\delta\widetilde{T}$ and final correction field $\delta T$

    \Statex
    \State \textbf{First stage: streamwise solve}
    \For{$j=1,\ldots,n_y$}
        \State Construct $\mathbf{r}^{(j)}$ using
        Algorithm~\ref{alg:apf_build_rhs}
        \State Construct $\mathbf{A}_x^{(j)}$ using
        Algorithm~\ref{alg:apf_build_x_matrix}
        \State Set the boundary entries of
        $\Delta t\,\mathbf{r}^{(j)}$ to zero
        \State Solve
        $\mathbf{A}_x^{(j)}
        \delta\widetilde{\mathbf{T}}^{(j)}
        =\Delta t\,\mathbf{r}^{(j)}$
        using the Thomas algorithm, as given in
        Equation~\eqref{eq:apf_x_matrix_system}
        \State Store the solution in the corresponding row of
        $\delta\widetilde{T}$
    \EndFor

    \Statex
    \State \textbf{Second stage: wall-normal solve}
    \For{$i=1,\ldots,n_x$}
        \State Construct $\mathbf{A}_y^{(i)}$ using
        Algorithm~\ref{alg:apf_build_y_matrix}
        \State Extract the corresponding column of
        $\delta\widetilde{T}$
        \State Solve
        $\mathbf{A}_y^{(i)}\delta\mathbf{T}^{(i)}
        =\delta\widetilde{\mathbf{T}}^{(i)}$
        using the Thomas algorithm, as given in
        Equation~\eqref{eq:apf_y_matrix_system}
        \State Store the solution in the corresponding column
        of $\delta T$
    \EndFor

    \Statex
    \State Update the temperature using
    Equation~\eqref{eq:apf_temperature_update}
    \State Apply the temperature boundary conditions
    \State Evaluate $R_{\mathrm{RMS}}$ using
    Equation~\eqref{eq:apf_convergence_residual}
    \State Append $n$ to the iteration history
    \State Append $R_{\mathrm{RMS}}$ to the residual history

    \If{$n\bmod N_s=0$}
        \State Store the current temperature field as a snapshot
        \State Record snapshot iteration $n$
    \EndIf

    \If{$R_{\mathrm{RMS}}\leq\varepsilon$}
        \State $\mathrm{converged}\gets\mathrm{true}$
        \State \textbf{break}
    \EndIf
\EndFor

\State Stop the timer and calculate computation time

\If{the final iteration is not already stored as a snapshot}
    \State Store the final temperature field and iteration
\EndIf

\If{not converged}
    \State Issue a warning that the tolerance was not reached
    within $N_{\max}$ iterations
\EndIf

\State \Return $T$, iteration history, residual history,
computation time, temperature snapshots, and snapshot iterations
\end{algorithmic}
\end{algorithm}

The APF method therefore advances the temperature field through
two sets of tridiagonal solves at each pseudo-time iteration.
The RMS residual measures convergence toward a solution of the
discrete steady energy equation. A small residual indicates
that the discrete equation is nearly satisfied, but does not
by itself establish spatial accuracy or grid independence.